% Glossary entries — use \gls{label} or \Gls{label} in text.
% Acronyms — use \acrfull{label}, \acrshort{label}, or \acrlong{label} in text.

%----------------------------------------------------------------------------------------
%	Terms
%----------------------------------------------------------------------------------------
\newglossaryentry{amnesiac}{
    name=amnesiac,
    description={OED:  (1913-) "One who is afflicted with amnesia." \footnote{  
    “Amnesiac, N.” Oxford English Dictionary, Oxford UP, December 2024, \url{https://doi.org/10.1093/OED/9207200444}.}}
}



\newglossaryentry{bsd}{
    name=BSD,
    description={Berkeley Software Distribution, a Unix-derived operating system and its variants (FreeBSD, OpenBSD, NetBSD)}
}

\newglossaryentry{coding}{
    name=coding,
    description={
    The OED described coding as (1946-)"3.  Computing. The conversion of information into code usable by a computer or other machine; the writing or editing of code for a computer program, application, etc. Also: computer code."
    \footnote{
    “Coding, N., Sense 3.” Oxford English Dictionary, Oxford UP, June 2025, \url{https://doi.org/10.1093/OED/4701717734}.
    }
    }
}


\newglossaryentry{configuration}{
    name=configuration,
    description={The OED describes configuration as (1962) "8. Computing. The way the constituent parts of a computer system are chosen or interconnected in order to suit it for a particular task or use; the units or devices required for this." 
    \footnote{
    “Configuration, N., Sense 8.” Oxford English Dictionary, Oxford UP, June 2026, \url{https://doi.org/10.1093/OED/1065280211}.}
    }
}


\newglossaryentry{immutable system}{
    name=immutable system,
    description={
    The OED describes immuatable as (1621) "1.b.  technical. Not subject to variation in different cases; invariable: used e.g. of markings which are the same in all the individuals of a species."  
    \footnote{“Immutable, Adj., Sense 1.b.” Oxford English Dictionary, Oxford UP, December 2025, \url{https://doi.org/10.1093/OED/1176087402}.}  An immutable system is one that cannot be changed.  In these computing systems, the intent of using an immutable system is to allow an operator to restore a system to its original state by power cycling it.
    }
}

\newglossaryentry{kernel}{
    name=kernel,
    description={The core component of an operating system that manages hardware resources and provides services to user-space programs}
}

\newglossaryentry{jail}{
    name=jail,
    description={A FreeBSD OS-level virtualisation mechanism that partitions the operating system into isolated environments}
}

\newglossaryentry{port}{
    name=port,
    description={A directory in the FreeBSD Ports Collection containing the files needed to automatically compile and install a third-party application}
}

\newglossaryentry{zfs}{
    name=ZFS,
    description={A file system used in FreeBSD adapted from Oracle.}
}

\newglossaryentry{poudriere}{
    name=Poudriere,
    description={Named after the French word for ``powderkeg,'' Poudriere is a package building program that will custom-compile programs from ports and configuration files.  A FreeBSD operator can then use their own tailored programs built to specifications while obtaining updates from the ports tree.}
}

\newglossaryentry{}{
    name=,
    description={}
}





%----------------------------------------------------------------------------------------
%	Acronyms
%----------------------------------------------------------------------------------------

\newacronym{cissp}{CISSP}{Certified Information Systems Security Professional}
\newacronym{csslp}{CSSLP}{Certified Secure Software Lifecycle Professional}
\newacronym{ccsp}{CCSP}{Certified Cloud Security Professional}
\newacronym{os}{OS}{Operating System}
\newacronym{ssh}{SSH}{Secure Shell}
\newacronym{tls}{TLS}{Transport Layer Security}
\newacronym{acl}{ACL}{Access Control List}
\newacronym{ids}{IDS}{Intrusion Detection System}
\newacronym{ips}{IPS}{Intrusion Prevention System}
\newacronym{vpn}{VPN}{Virtual Private Network}
\newacronym{ram}{RAM}{Random Access Memory}
\newacronym{vnet}{VNET}{Virtual Network}
\newacronym{gb}{GB}{GigaBytes}
\newacronym{dvd}{DVD}{Digital Video Disc}
\newacronym{rw}{RW}{Read Write}
\newacronym{lto}{LTO}{Linear Tape Open}
\newacronym{ssd}{SSD}{Solid State Drive}
\newacronym{sas}{SAS}{Serial Attached SCSI}
\newacronym{sata}{SATA}{Serial AT Attachment}
\newacronym{nic}{NIC}{Network Interface Card}
\newacronym{gpu}{GPU}{Graphics Processing Unit}
\newacronym{nist}{NIST}{National Institute of Standards and Technology}
